\section{Artinian Rings}

\begin{exercise}\label{ex8.1}
	Use \Cref{7.16} on local ring $A_{\mathfrak{p}_i}$. Then apply the contraction. If $\mathfrak{q}_i$ is isolated, then $\mathfrak{p}_i$ is minimal.
\end{exercise}

\begin{exercise}\label{ex8.2}
	\begin{enumerate}
		\item[(i) $\Rightarrow$ (ii)] Use \Cref{8.1}.
		\item[(iii) $\Rightarrow$ (i)] Use \Cref{8.5}.
	\end{enumerate}
\end{exercise}

\begin{exercise}\label{ex8.3}
	\begin{enumerate}
		\item[(i) $\Rightarrow$ (ii)] By \Cref{8.7}, we can assume that $A$ is an Artinian local ring. By \Cref {8.4}, $\mathfrak{m}^n=0$ for some $n$, so we have a chain $A=\mathfrak{m}^0\supseteq\mathfrak{m}^1\supseteq\cdots\supseteq\mathfrak{m}^n=0$.\par
		By \hyperref[7.9]{Zariski's lemma}, $A/\mathfrak{m}$ is a finitely generated $k$-module. By \Cref{8.5}, $A$ is Noetherian, so $\mathfrak{m}^i$ is a finitely generated $A$-module,  and therefore $\mathfrak{m}^i/\mathfrak{m}^{i+1}$ is a finitely generated $A/\mathfrak{m}$-module. Use
		\begin{align*}
			&\dim_k\left(\mathfrak{m}^i\middle/\mathfrak{m}^{i+1}\right)=\dim_{A/\mathfrak{m}}\left(\mathfrak{m}^i\middle/\mathfrak{m}^{i+1}\right)\dim_k\left(A\middle/\mathfrak{m}\right);&&\\
			&\dim_k\left(\mathfrak{m}^i\right)=\dim_k\left(\mathfrak{m}^{i+1}\right)+\dim_k\left(\mathfrak{m}^i\middle/\mathfrak{m}^{i+1}\right),&&\text{where }0\rightarrow\mathfrak{m}^{i+1}\rightarrow\mathfrak{m}^i\rightarrow\mathfrak{m}^i/\mathfrak{m}^{i+1}\rightarrow0
		\end{align*}
		to finish the proof.
		\item[(ii) $\Rightarrow$ (i)] $A$ is a $k$-vector space of finite dimension. Use \Cref{6.10}.
	\end{enumerate}
\end{exercise}

\begin{exercise}\label{ex8.4}
	\begin{enumerate}
		\item[(i) $\Rightarrow$ (iii)] The finitely-generated property is preserved for localization and quotient of modules.
		\item[(ii) $\Rightarrow$ (iii)] $B_\mathfrak{p}/\mathfrak{p}B_\mathfrak{p}$ is a finitely generated $k(\mathfrak{p})$-algebra. By \hyperref[7.5]{Hilbert's basis theorem}, any finitely generated $k(\mathfrak{p})$-algebra is Noetherian. Then use \Cref{ex8.2} and \Cref{ex8.3}.
		\item[(iii) $\Rightarrow$ (ii),(iv)] $B_\mathfrak{p}/\mathfrak{p}B_\mathfrak{p}$ is Noetherian and integral over the field $k(\mathfrak{p})$. If $\mathfrak{q}\in B_\mathfrak{p}/\mathfrak{p}B_\mathfrak{p}$ is prime, $\left(B_\mathfrak{p}/\mathfrak{p}B_\mathfrak{p}\right)/\mathfrak{q}$ is an integral domain, so it is a field by \Cref{5.7}. Hence, every prime ideal in $B_\mathfrak{p}/\mathfrak{p}B_\mathfrak{p}$ is maximal, which means $\operatorname{dim}\left(B_\mathfrak{p}/\mathfrak{p}B_\mathfrak{p}\right)=0$. Thus, $B_\mathfrak{p}/\mathfrak{p}B_\mathfrak{p}$ is Artinian by \Cref{8.5}. Then use \Cref{8.3}. 
	\end{enumerate}
	A counter example is $f:\mathbb{Q}\hookrightarrow\overline{\mathbb{Q}}$ where $\overline{\mathbb{Q}}$ is an algebraic closure of $\mathbb{Q}$.
\end{exercise}

\begin{exercise}\label{ex8.5}
	Use \hyperref[noethernormalization]{Noether normalization lemma} and let $\mathfrak{m}$ be a maximal ideal in $k\left[y_1,\cdots,y_r\right]$. By \Cref{ex8.4}, $A_\mathfrak{m}/\mathfrak{m}A_\mathfrak{m}$ is finite over $k\left[y_1,\cdots,y_r\right]/\mathfrak{m}=k$, so by \Cref{ex8.3} and \Cref{8.3}, $A_\mathfrak{m}/\mathfrak{m}A_\mathfrak{m}$ only has finite many maximal ideals, correspond to finite many points in $X$.
	
	Let $N$ be the number of generators of $A$ over $k\left[y_1,\cdots,y_r\right]$. The $k$-dimension of $k$-vector space $A_\mathfrak{m}/\mathfrak{m}A_\mathfrak{m}$ is uniformly bounded by $N$. If $\mathfrak{n}_1,\cdots\mathfrak{n}_M$ are maximal ideals in $A_\mathfrak{m}/\mathfrak{m}A_\mathfrak{m}$, by \Cref{1.10}, $A_\mathfrak{m}/\mathfrak{m}A_\mathfrak{m}\rightarrow\prod_{i=1}^M\left.\left(A_\mathfrak{m}/\mathfrak{m}A_\mathfrak{m}\right)\middle/\mathfrak{n}_i\right.$ is surjective, so $\dim_k\left(A_\mathfrak{m}/\mathfrak{m}A_\mathfrak{m}\right)\geqslant\sum_{i=1}^M\dim_k\left(\left.\left(A_\mathfrak{m}/\mathfrak{m}A_\mathfrak{m}\right)\middle/\mathfrak{n}_i\right.\right)\geqslant M$.
\end{exercise}

\begin{exercise}\label{ex8.6}
 	Consider $(A/\mathfrak{q})_{\overline{\mathfrak{p}}}$, which is an Artinian local ring. Use the following claims then \Cref{6.7}.
 	\begin{itemize}
 		\item \textit{If $A$ is an Artinian local ring, $\mathfrak{m}$ its only maximal ideal. Then every ideal in $A$ is $\mathfrak{m}$-primary.}\\ Radical is the intersection of the only prime ideal.
 		\item \textit{There is a 1-1 correspondence between $\mathfrak{p}$-primary ideals in $A$ and $S^{-1}\mathfrak{p}$-primary ideals in $S^{-1}A$.}\\
 		Use \Cref{3.11.1} and \Cref{4.9}.
 	\end{itemize}
\end{exercise}